% SEGMENT OLP-0494-S01
% part: intuitionistic-logic
% chapter: introduction
% section: syntax

\documentclass[../../../include/open-logic-section]{subfiles}

\begin{document}

\olfileid{int}{int}{syn}

\olsection{உள்ளுணர்வுவாதத் தருக்கத்தின் தொடரியல்}

உள்ளுணர்வுவாதத் தருக்கத்தின் தொடரியல், கூற்றுக்கணிப்புத் தருக்கத்தின்
தொடரியலைப் போன்றதே. செவ்வியல் கூற்றுக்கணிப்புத் தருக்கத்தில் சில
இணைப்பிகளைப் பிற இணைப்பிகளால் வரையறுக்க முடியும்; எடுத்துக்காட்டாக,
$!A \lif !B$ ஐ $\lnot !A \lor !B$ ஆலும், $!A \lor !B$ ஐ
$\lnot(\lnot !A \land \lnot !B)$ ஆலும் வரையறுக்கலாம். எனவே செவ்வியல்
தருக்கத்தின் விளக்கங்கள் சில இணைப்பிகளை இவ்வரையறைகளுக்கான சுருக்கங்களாக
அடிக்கடி அறிமுகப்படுத்துகின்றன. இரண்டு விதிவிலக்குகளைத் தவிர,
உள்ளுணர்வுவாதத் தருக்கத்தில் இவ்வாறு செய்ய முடியாது: $\lnot !A$ ஐ
$!A \lif \lfalse$ என்பதன் சுருக்கமாக வரையறுக்கலாம்---அடிக்கடி அவ்வாறே
வரையறுக்கப்படுகிறது. அப்போது, $\lfalse$ ஐயே வரையறுக்கக் கூடாது!{} மேலும்,
செவ்வியல் தருக்கத்தைப் போலவே, $!A \liff !B$ ஐ $(!A \lif !B) \land
(!B \lif !A)$ என வரையறுக்கலாம்.

கூற்றுக்கணிப்பு உள்ளுணர்வுவாதத் தருக்கத்தின் !!^{formula}s,
\emph{!!{propositional variable}s}, கூற்று மாறிலி~$\lfalse$ மற்றும்
\emph{தருக்க இணைப்பிகள்} ஆகியவற்றிலிருந்து கட்டப்படுகின்றன. நம்மிடம்
பின்வருவன உள்ளன:

\begin{enumerate}
\item !!{propositional variable}s $\Obj p_0$, $\Obj p_1$, \dots இன்
  ஓர் !!^a{denumerable} கணம்~$\PVar$.
\item !!{falsity}க்கான கூற்று மாறிலி~$\lfalse$.
\item தருக்க இணைப்பிகள்: $\land$ (இணைப்பு), $\lor$ (பிரிப்பு), $\lif$
  (நிபந்தனைவாக்கம்).
\item நிறுத்தற்குறிகள்: (, ), மற்றும் காற்புள்ளி.
\end{enumerate}

\begin{defn}[வாய்பாடு]
\ollabel{defn:formulas} கூற்றுக்கணிப்பு உள்ளுணர்வுவாதத் தருக்கத்தின்
\emph{!!{formula}s} இன் கணம்~$\Frm[L_0]$ பின்வருமாறு தொகுமுறையில்
வரையறுக்கப்படுகிறது:
\begin{enumerate}
\item $\lfalse$ ஓர் அணு !!{formula}.

\item ஒவ்வொரு !!{propositional variable}~$\Obj p_i$ உம் ஓர் அணு
  !!{formula}.

\item $!A$, $!B$ ஆகியவை !!{formula}s எனில், $(!A \land !B)$ ஓர்
  !!a{formula}.

\item $!A$, $!B$ ஆகியவை !!{formula}s எனில், $(!A \lor !B)$ ஓர்
  !!a{formula}.

\item $!A$, $!B$ ஆகியவை !!{formula}s எனில், $(!A \lif !B)$ ஓர்
  !!a{formula}.

\tagitem{limitClause}{வேறு எதுவும் ஓர் !!a{formula} அல்ல.}{}
\end{enumerate}
\end{defn}

மேலே அறிமுகப்படுத்திய முதன்மை இணைப்பிகளுடன், பின்வரும்
\emph{வரையறுக்கப்பட்ட} குறிகளையும் பயன்படுத்துகிறோம்: $\lnot$ (மறுப்பு),
$\liff$ (!!{biconditional}). வரையறுக்கப்பட்ட செயற்குறிகளைக் கொண்டு
கட்டப்பட்ட வாய்பாடுகள் பின்வருமாறு புரிந்துகொள்ளப்பட வேண்டும்:
\begin{enumerate}
\item $\lnot !A$ என்பது $!A \lif \lfalse$ என்பதன் சுருக்கம்.

\item $!A \liff !B$ என்பது $(!A \lif !B) \land (!B \lif !A)$
என்பதன் சுருக்கம்.
\end{enumerate}

$\lnot$ முறைப்படி ஒரு சுருக்கமாகக் கருதப்பட்டாலும், அது முதன்மையானது
என்பதுபோல் சில நேரங்களில்~$\lnot$ க்கென வெளிப்படையான விதிகளையும் வரையறைக்
கூறுகளையும் தருவோம். இது பெரும்பாலும் பயிற்சி வினாக்களை வகுக்க உதவுகிறது.

\end{document}
